\documentclass{article}
\usepackage{graphicx}
\usepackage{subcaption}
\usepackage[colorlinks=true, pdfstartview=FitV, linkcolor=blue, citecolor=blue, urlcolor=blue]{hyperref}
\usepackage[authoryear]{natbib}
\usepackage[affil-it]{authblk}
\usepackage{fancyhdr}
\usepackage{pdflscape}
\usepackage{longtable}
\usepackage{float}
\usepackage{booktabs}

\renewcommand\Affilfont{\fontsize{9}{10.8}\itshape}

\bibliographystyle{apalike} 

\begin{document}
\title{Status advantage without a Matthew effect: cross-cultural evidence from prosocial reputation updating in economic games} 

\pagestyle{fancy}

\maketitle
[Note for the future: not sure what to say about the cross-culturalness because I don't believe in culture, or about why small villages, I leave it to my beloved ethnographers colleagues to make me suggestions]
\bigbreak
To explain the puzzling human amount of cooperation among humans,
economists and behavioral scientiss have long proposed reputation as a central mechanism [Wu, J., Balliet, D. \& Van Lang, other refs from Frey 11-17]. Reputation, the reasoning goes, guides people to choose whether to cooperate [REF], with whom to cooperate [REF]. and in turn incentivizes `good' behaviour, as a way to improve one's reputation. This literature has been extremely succesful in showing why explaining why reputation `works': how it reveals cooperative tendencies, how it benefits one to have a good reputation, and how it in turn stabilizes cooperation. 
\bigbreak
The literature on cooperation and reputation has, however, taken for granted that reputation is accurate [REF Milinsky], or at least a noisy but unbiased proxy [Julien Proc B, REF Hilbe 2018 noisy]. It has, often, assumed so by design, by operationalizing reputation as 'image scoring', that is, as behavioural history. In canonical models (Nowak Sigmund) as in experiments, where 'reputation' is implemented by merely showing the past play (Milinsky Nature and Science). This axiom that reputation is accurate is rarely tested [ref four puzzles], and when it is, the interest is in the degree of agreement, rather than disagreement. 
\bigbreak
Relatedly, reputation has mostly been treated as an independant variable to be manipulated, and the focus has been on its consequences, rather than reputation values themselves. Models needs to assume that reputation varies within a population, as a reflect of intrinsic variation in altruistic tendencies -- if everyone is similar, there is no need to track anything, but little attention has been given to the distribution of reputation values in populations.
\bigbreak
Sociologists, for their part, have taken a radically opposite view of reputation. They analyse reputation as an aspect of inequality: reputation is a form of status, a "dimension of inequality with its own sustaining social dynamic" (AJS 2013, p. 4) (REF Weber). Therefore, sociologists have rather looked into how reputation could be 'socially contstructed' [REF Lynn], that is, persistent [REF Tilly], self-perpetuating [REF Ridgeway] and arbitrary [REF Frey]. Concretely, some actors initially, can obtain a higher reputation due to chance or having more resources, then some process locks in these inequality, in what Merton dubbed a 'Matthew effect' [REF Merton] - also later called 'cumulative advantage', meaning that current advantage causes later stronger advantage (Ref di prete). Multiple examples of such 'success breed success' [REF vanderijt] dynamics have been found, in multiple domains [REF Merton, vanderijt, sagalnik]. 


To justify cumulative advantage dynamics, Merton proposed that a same signal might be interpreted differently based on the sender's status: as Ridgeway put it, "the same idea ‘sounds better’ coming from the advantaged than the disadvantaged" (Ridgeway address). This hypothesis, `status advantage', has often been taken as equivalent to a Matthew effect -- even though formal models show that status advantage amplify reputational inequality, but not necessarily decouples it from underlying qualities - what Lynn [REF] called 'weak' and 'strong' social construction.

In this paper, we analyse the dynamics of reputation under the light of this sociological insight. In eight field sites, we measured reputation before and after a dictator game - a common measure of 

This address a double literature gap. First, in the cooperation literature, we test whether reputation is simply a mirror of behaviour, or a self-sustaining dimension of inequality [Ridgeway]. Second, it extends the sociological perspective, to a new, arguably more realistic domain. Most evidence in sociology is based on domains where reputation is quite different from prosociality reputation: extremely skewed, extreme quality uncertainty and impossible to know everyone. This makes the dynamics probably closer to capital accumulation than to beliefs updating.


Here, we import this insight from sociology, and test whether it works for prosociality reputation. It is quite clear that Apple, Taylor Swift and Daron Acemoglu benefit from status advantage, but does a generous Jane Doe does too?

Literature gaps:
- cumulative advantage and status advantage tested at the same time. But to do it clean, you should hold the signal constant?






Fiske: status predicts competence, but not warmth! So maybe prosocial reputation is less socially biased?






Di Prete: strict (Mertonian) vs loose (Blau-Duncan) version of cumulative advantage.





- does it applies to prosociality reputation? To paraphrase Ridgeway, does an action also seems less stingy when it comes from an established person?
- ESF framework is based on competence/performance evaluation



\section{Model}
At time 1, an agent posess a reputation $r_1 > 0$, which is an estimator of its secrete generosity $g$. She sends a signal $s > 0$, which feeds into his reputation this way:
$$r_2 = (1-\lambda) r_1 + \lambda s
$$

If $s=g$, or if $s$ is a noisy signal centered on $g$, then 

Now assume that there is status advantage: signal reception is biased by existing reputation. Formally, if one's actual generosity is $g$, the received signal is $g + \beta s$, with $\beta$ the status advantage term. 

In other words,
$$
r_2 = (1-\lambda) r_1 + \lambda(g+\beta r_1) = (1-\lambda(1 - \beta))r_1 + \lambda g
$$

Long term reputation is given by the fixed point of the equation: 
$$
x = (1-\lambda(1 - \beta))x + \lambda g \Longleftrightarrow  x= \frac g {1-\beta}
$$

question: is it capital accumulation, or belief updating?
\section{Data and methods}
Pre/post reputational evaluations, dictator game in between, revelation or not (manipulation). Also, we ask whether the revealed decisions were "appropriate".

We do a diff and diff model, since 

\section{Results}



\section{Discussion}




Limitations: is the status/reputation analogy valid at all? Ivan's objection: if reputation is viewed as how many people know me, then our measure speak little about this.

\end{document}